\cotefile{93}
\scanpage{221}
% top line and annotated display
\deleted{En général}, On trouve alors un \unsure{homom.|morph.|extension}
\[
\underset{\substack{\vert\\\text{faisceau }(\deleted{\text{Zar}})\\
 \text{des sections}\\\text{de }G}}{\underline{\Gamma}}(G)/\underline E
\underset{\deleted{\unsure{\log_G}}}{\longrightarrow}
Z^1\underbrace{C_{S/R}^{\bullet}(\underline g^{\unsure{*}})}_{
 \substack{\text{complexe des}\\\text{diff. rel/R à}\\\text{coefft en }\underline g}}
\,/\,d^0(E)
\]

% middle, beside the underlined proposition
D'autre part, d'après la théorie des Modules stratifiés, \uline{tout op. différentiel} \uline{``à coeff cts''} $D:\underline g\longrightarrow \unsure{\mathcal M|M}$
% slanted, circled
\margin{spécial à car. $0$ ?}
\deleted{définit un op. diff} \uline{se factorise de façon canonique en}
\diagram{\begin{tikzcd}[ampersand replacement=\&,column sep=large]
\underline g \arrow[r,"{c=d^0}"] \arrow[rr,bend right=25,"D"'] \&
 Z^1C_{S/R}(\underline g) \arrow[r,"{\widetilde D}"] \& \mathcal M
\end{tikzcd}}

% parenthesis below the first diagram
($\widetilde D$ est loc. induit par un op. diff. sur $\underline g\otimes\Omega^1_{S/R}$ ; \deleted{\unsure{de tels}} \deleted{tous les} $\Delta$, \deleted{\unsure{donnant}} \inserted{avec} $\Delta\circ d^0=D$, ont $\overline{\mathrm m}$ restriction à $Z^1C_{\unsure{c|S}/R}(\underline g)$..)

% closing diagram
Donc on trouve un composé
\diagram{\begin{tikzcd}[ampersand replacement=\&,column sep=large]
\underline{\Gamma}(G) \arrow[r,"{g^{-1}dg}"] \arrow[rr,bend right=25,"{\widetilde D}"{draw,circle,inner sep=1pt}] \&
 Z^1C_{S/R}(g) \arrow[r,"{\widetilde D}"] \& M
\end{tikzcd}}
% archive sheet number: 220

\scanpage{222}
% opening system of data, lines 1–4
1) Soit $f:X\longrightarrow S$ morphisme lisse d'espaces analytiques (lisses)

$\omega$ une 1-forme fermée \inserted{(holomorphe)} relative sur $X/S$

$a$ et $b$ deux sections \inserted{holomorphes} de $X$ sur $S$, telles que $\forall s\in S$, $a(s)$ et $b(s)$ \inserted{soient dans la $\overline{\mathrm m}$ comp. connexe de $X_s$}

% local paths, lines 5–7
Localement sur $S$, on peut trouver un arc \deleted{\unsure{diff.|diff}} \inserted{(différentiable par morceaux)} $C_s:[0,1]=I\longrightarrow X_s$, joignant $a(s)$ et $b(s)$, et dépendant \unsure{continument|continûment} de $s$.

% primitive and the first integral, lines 8–12
Utilisant le fait que $\omega$ est fermée, donc \unsure{localement|localem\textsuperscript{t}} de la forme $d_{X/S}\varphi$, on constate que la fonction sur $S$
\[
 I(s)=\int_{\deleted{\unsure{a(s)}}C_s}^{\deleted{\unsure{b(s)}}}\omega_s
       =\int_{a(s)}^{b(s)}\omega_s
\]
est holomorphe en $s$. [\unsure{En|On|En effet} \inserted{\unsure{se plaçant|plus précisément} au voisinage de $s_0\in S$ et \unsure{se|en}} \unsure{ramenant|ramène|supposons} au cas où \deleted{\unsure{$a(s_0)$ et|$a(s_0)$}} $a(s_0)$ et $b(s_0)$ sont assez voisins pour que dans un ouvert $U$ contenant $a(s_0),b(s_0)$, on ait $\omega|U=d\varphi$, donc pour $s$ voisin de $s_0$, $I(s)=\varphi b(s)-\varphi a(s)$].

% change of paths and equation (1)
Changeant le système des $C_s$ revient à les remplacer par $C'_s=C_s+Z_s$, $Z_s$ un 1-cycle, donc on trouve
\[
 I'(s)=I(s)+\underbrace{\int_{Z_s}\omega_s}_{\deleted{\unsure{Z}}}
 \tag*{(1)}
\]

% bottom and equation (2)
Considérons une expression formelle (avec des $\omega_i$ comme $\omega$, et des op. diff. \inserted{(hol.)} $D_i$ sur $S$ ]
\[
 \sum D_{i\,s}\int_{a(s)}^{b(s)}\omega_{i\,s}
 \tag*{(2)}
\]
% archive sheet number: 221

\scanpage{223}
% opening, (3)
Cette fonction est déterminée, indépendamment du choix de $C_s$, modulo une expression de la forme
\[
 \sum D_{i\,s}\int_{Z_s}\omega_{i\,s}
 \deleted{=\int_{Z_s}\sum_i D_{i\,s}\omega_{i\,s}}
 \tag*{(3)}
\]

% cancelled NB and the definition
\deleted{(\uline{NB} Le fait que $Z_s$ est un 1-\uline{cycle} permet \unsure{de}}

Supposons que les $D_i$, $\omega_i$ soient tels que l'expression précédente soit nulle, quelle que soit \deleted{\unsure{est}} la famille continue $s\unsure{\longmapsto|\longrightarrow}Z_s$ de 1-cycles sur un ouvert de $S$. \deleted{\unsure{$Z$|$\Sigma$}} \deleted{Alors l'expression} \inserted{On dira} que \deleted{les} $\mu=(D_i,\omega_i)_{i\in I}$ \uline{définit une équation de Picard-Fuchs sur} $X/S$. Alors pour deux sections hol. quelc. $a,b$ de $X/S$, l'expression
\[
 \mu(a,b)=\sum_i D_{i\,s}\int_{a(s)}^{b(s)}\omega_{i\,s}
 \tag*{(4}
\]
est bien déterminée, et est une fonction holomorphe sur $S$. Bien sûr on aura
\[
 \mu(a,b)+\mu(b,c)=\mu(a,c).\tag*{(5)}
\]

% last lines, continuing on the next sheet
Cherchons une condition suffisante maniable, \inserted{sous forme différentielle} \deleted{\unsure{pour}} \deleted{un \unsure{critère} \uline{infinitésimal}, et non plus global,} pour que $\mu$ soit de Pic. Fuchs. Notons
% archive sheet number: 222

\scanpage{224}
% continuation and framed definition
que \deleted{\unsure{pour}} sur
\[
 \boxed{\Phi=f_*(Z\underline{\Omega}^{1}_{X/S})/d f_*(\underline{\mathcal O}_X)}
 \tag*{(6)}
\]
, il y a une connexion intégrable naturelle, grâce à la formule de Cartan --- Les $\omega_i$ définissent des sections $\omega'_i$ de $\Phi$, et l'on peut former les $D_i\deleted{\unsure{\omega}}\omega'_i$. Je dis qu'il suffit qu'on ait
% slanted
\margin{attention, ne marche tel quel que si $X\longrightarrow S$ est steinien\unsure{!|?}}
\[
 \sum D_i\omega'_i=0.\tag*{(7)}
\]

% calculation using a connection
Lorsque $X/S$ est muni d'une connexion intégrable, alors le calcul de $\sum D_i\omega'_i$ se fait ainsi : \inserted{$f_*(\Omega^*_{X/S})$ sera aussi stratifié, donc} chaque $D_i$ se remonte \deleted{par la stratification en} \inserted{en un opérateur} $\overline D_i$ \deleted{sur} $\underline{\Omega}^*_X$ \inserted{commutant à la diff. ext. $d_{X/S}$} et \deleted{\unsure{comme} $\omega_i$ est une forme fermée, on peut former}

% entire enclosed block cancelled by several large oblique crosses
\deleted{\[
\boxed{\begin{aligned}
 (8)\quad &\alpha_{\deleted{\unsure{D_i}}\overline D_i}\omega_i
 =\alpha_{\overline D_i^0}\omega_i+c_i\,\omega_i\\
 &\left\{\begin{aligned}
 \overline D_i&=\overline D_i^0+c_i\,\mathrm{id}
 &\left\{\begin{aligned}&\overline D_i^0\text{ à coeff. cts}\\
                         &c_i\text{ section de }\underline{\mathcal O}_X\end{aligned}\right.\\
 (9)\quad\alpha_{\overline D_i^0}\,d\varphi&=\deleted{\unsure{\alpha_i}}\overline D_i^0\varphi
 \end{aligned}\right.
\end{aligned}}
\]}

% rewritten (8) at the foot
Considérons alors
\[
 \sum \deleted{\unsure{\alpha}}(\overline D_i)\omega_i=\overline\omega
 \tag*{(\deleted{\unsure{9}}8)}
\]
qui est encore une forme fermée. On vérifie alors \inserted{aussitôt} que l'on a
% archive sheet number: 223

\scanpage{225}
% (9) and (10)
\[
 \sum D_i\omega^{\unsure{'}}_i=\overline\omega'\qquad\text{dans }\Phi
 \tag*{(9)}
\]
et la condition (7) signifie que l'on a, localement au dessus de $S$, une relation
\[
 \overline\omega\; \bigl(=\sum\overline D_{\deleted{\unsure{s}}i}\,\omega_i\bigr)=d\varphi,\qquad
 \varphi\in\Gamma(X,\mathcal O_X).\tag*{(10)}
\]

% period over a cycle
Pour voir que cette relation implique bien la nullité de (3), il suffit de prouver
\[
 \sum_i D_{i\,s}\int_{Z_s}\omega_{i\,s}=\int_{Z_s}\overline\omega
 \tag*{(11)}
\]
qui implique, si $\overline\omega=d\varphi$, que le premier mb est égal à $\int_{\partial Z_s}\varphi=0$ car $\partial Z_s=0$.

% general relation and its hypotheses
En fait, (11) provient de la relation générale
\[
 D_s\int_{Z_s}\omega_s=\int_{Z_s}\overline D_s\,\omega_s
 \tag*{(12)}
\]
valable pour tt op. diff. hol. $D$ sur $S$, et tte forme \inserted{holomorphe fermée} $\omega_{\deleted{\unsure{s}}}$ sur $X$ [le fait que $\omega$ est fermée implique que $\int_{Z_s}\omega$ est fonction holomorphe en $s$, et donne un sens au premier mb de (12)]. Pour prouver cette
% archive sheet number: 224

\scanpage{226}
% continuation; the bracketed draft is crossed through with a wavy stroke
dernière, on \deleted{[peut se ramener, en décomposant $Z_{s_0}$ en \deleted{\unsure{petits arcs}} \inserted{petits} chemins $a_i^0\longrightarrow a^0_{i+1}$, et transportant \deleted{par} horizontalement les $a_i^0$, \deleted{\unsure{à}}]} va prouver la relation plus générale
\[
 D_s\int_{a(s)}^{b(s)}\omega_s=\int_{a(s)}^{b(s)}(\overline D\omega)_s,
 \tag*{(13)}
\]
valable pour des \unsure{limites|éléments|chemins} d'intégration $a(s)$, $b(s)$ qui sont des \uline{sections} \uline{horizontales} pour la connexion. [N.B. Bien \deleted{entendu}, on suppose, pour donner un sens à cette formule, qu'on choisit les $\overline{\mathrm m}^{\mathrm s}$ $C_s$ au deux mbs, donc (13) s'écrit avec tte précision
\[
 D_s\int_{C_s}\omega_s=\int_{C_s}(\overline D\omega)_s\qquad].
 \tag*{(13 bis)}
\]

% subdivision, ending in the middle of a sentence
Pour prouver (13), on est ramené, en \inserted{se plaçant au voisinage de $s_0\in S$ et} décomposant le \unsure{domaine|chemin} $\int_{a(s)}^{b(s)}$, au cas où $a(s_0),b(s_0)$ sont assez voisins, de sorte qu'on peut se placer dans un ouvert $U$ où l'on a $\omega=d\varphi$. Mais alors le $1^{\mathrm{er}}$ mb est $D_s\bigl(\varphi(b(s))-\varphi(a(s))\bigr)$, le
% archive sheet number: 225

\scanpage{227}
% continuation of the proof of (13)
deuxième est $\int_{a(s)}^{b(s)}\overline D\,d\varphi
=\int_{a(s)}^{b(s)}d\,\overline D\varphi
=(\overline D\varphi\unsure{]|)}(b(s))-(\overline D\varphi)(a(s))$,
\deleted{et les} et les deux expressions sont égales, grâce au fait que $a,b$ sont des sections horizontales. Notons aussi que ceci montre, si $a,b$ sont des sections \uline{horizontales}, qu'on aura
\[
 \mu(a,b)=\int_{a(s)}^{b(s)}\bigl(\sum\overline D_i\omega_i\bigr)_s
 =\varphi(b(s))-\varphi(a(s)),\tag*{(14)}
\]
où $\varphi$ est défini par la relation (10).

% sufficiency in general
Prouvons la suffisance de (7) pour la nullité de (3) en général. On va prouver ``plus généralement'' qu'on aura alors
\[
 \sum D_{i\,s}\int_{C_s}\omega_{i\,s}=0
\]
lorsque $C_s$ est un \unsure{1-cycle|1-cycles|1-chaîne} dans $X_s$, variant différentiablement avec $s$, dont le bord est ``horizontal'' en $s$, ou ce qui revient au $\overline{\mathrm m}$
\[
 \sum_i D_{i\,s}\int_{a(s)}^{b(s)}\omega_{i\,s}=0
\]
si $a,b$ sont des sections \uline{horizontales} de $X$ sur \deleted{\unsure{$S$}} \inserted{un ouvert de $S$.} \deleted{Pour} Se plaçant au voisinage d'un
% archive sheet number: 226

\scanpage{228}
% continuation, reduction about a point
point $s_0\in S$ on est ramené \unsure{comme|ainsi|encore} au cas où $a(s_0)$ est proche de $b(s_0)$, de sorte qu'on peut \deleted{\unsure{se pla}} se placer dans un ouvert de $X$ assez petit pour qu'il y ait une \unsure{connection|connexion}. Et on gagne.

% alternate formulation
On peut aussi noter, \deleted{\unsure{qu'on|quitte}} pour prouver la \deleted{formule} nullité de (3) comme conséquence de (7), que l'on a
\[
 \sum D_{i\,s}\int_{Z_s}\omega_i
 =\int_{Z_s}\bigl(\sum\overline D_i\,\omega'_i\bigr)_s
 \tag*{\deleted{\unsure{(15)}}}
\]
où le deuxième mb a un sens, car $\int_{Z_s}\overline\omega_s=0$ si $\overline\omega_s$ est un bord, \unsure{or|et} $\bigl(\sum\overline D_i\,\omega'_i\bigr)_s$ est une forme définie modulo un bord. La formule envisagée \deleted{\unsure{équivaut à}} \inserted{revient à}
\[
 \boxed{D_s\int_{Z_s}\omega'_s=\int_{Z_s}(\overset{\deleted{\unsure{-}}}{D}\omega')_s}
 \tag*{(15)}
\]
qui donne la propriété transcendante essentielle de la connexion intégrable canonique sur $\Phi$. Nous n'allons pas prouver (15) ici, et l'admettrons, mais notons que
% archive sheet number: 227

\scanpage{229}
% consequence of (15)
(15) nous donne une \uline{réciproque} pour la condition envisagée pour que $(D_i,\unsure{\varphi_i|\omega_i})$ soit de Picard-Fuchs : il f. et s. que l'on ait (7). Cela résulte du th. de De Rham sur les fibres $X_s$ .....

% heading and computation with horizontal sections
\uline{Calcul de} $\mu(a,b)$ par voie ``algébrique'', sans calcul d'intégrale. \inserted{C'est local sur $S$.} La formule (5) nous ramène à calculer $\mu(a,b)$ \inserted{au voisinage de $s_0$,} pour \deleted{$a,b$} $a(s_0)$ et $b(s_0)$ assez voisins. Cela nous permet alors de supposer que sur $X/S$ il y a une connexion intégrable. On peut $\overline{\mathrm m}$ supposer, par un choix convenable, que les sections $a,b$ \unsure{soient|sont} horizontales. Alors on trouve (14)
\[
 \mu(a,b)=\varphi(b(s))-\varphi(a(s))
\]
où $\varphi$ est défini par (10).

% beginning of the nonhorizontal computation
Il y aurait lieu d'expliciter aussi une \unsure{forme|formule} explicite, sans supposer que $a(s)$ et $b(s)$ soient horizontales. Faisons le calcul au voisinage de $s_0$, et \unsure{soit|soient} $a_0,b_0$ les
% archive sheet number: 228

\scanpage{230}
% horizontal sections; three displays
sections horizontales au voisinage de $s_0$, telles que
\[
 a_0(s_0)=a(s_0),\quad b_0(s_0)=b(s_0).
\]
On a
\[
 \mu(a,b)=\underbrace{\mu(a_0,b_0)}_{
  \deleted{\int_{a_0}^{\unsure{b_0|b}}\varphi}\ \varphi(b_0)-\varphi(a_0)}
 +\mu(b_0,b)-\mu(a_0,a)
\]
donc
\[
 \mu(a,b)(s_0)=\bigl[\varphi(b(s_0))-\varphi(a(s_0))\bigr]
 +\left(\sum_i D_{i\,s}\int_{b_0(s)}^{b(s)}\omega_{i\,s}
 -\sum_i D_{i\,s}\int_{a_0(s)}^{a(s)}\omega_{i\,s}\right)_{s=s_0}
\]
Il faut calculer les deux termes de \unsure{rabiot}, c'est une histoire purement locale.

% lower half cancelled by large diagonal strokes
\deleted{Notons qu'on a, grâce à la connexion de $X/S$
\[
 da\in\Gamma(S,a^*(\unsure{\mathscr T|\mathcal T|\mathscr V}_{X/S})\otimes\Omega^1_S)
\]
d'où, en contractant avec
\[
 a^*(\omega)\in\Gamma(S,\underbrace{a^*(\Omega^1_{X/S})}_{
 \substack{\text{dual de}\\a^*(\mathscr T_{X/S})}})
\]
% slanted, outside the diagonal cancellation on the manuscript
}
\margin{donner une formule..... !}
\deleted{une section, que nous noterons $a^*(\omega)$
\[
 a^*(\omega)\in\Gamma(S,\Omega^1_S)
\]
Cette \deleted{section} \inserted{forme} est encore \uline{fermée}, \inserted{$\omega$ l'étant} [N.B. pour le voir définir $a^*(\omega)$ pour tte forme (degré quelconque, fermée ou non) sur $X$, et noter que $a^*$ commute à la différentielle ext.)}
% archive sheet number: 229

\scanpage{231}
% relative scheme and the sequence in 1)
\diagram{\begin{tikzcd}[ampersand replacement=\&,row sep=small]
 S\arrow[d]\&\text{\unsure{schéma} relatif}\\T\&
\end{tikzcd}}
1) Soit $F\xrightarrow{D}G\xrightarrow{D'}H$ une suite d'op. différentiels, \inserted{d'ordre 1} qui est ``formellement exacte''. Alors pour tt $P$, la suite
\diagram{\begin{tikzcd}[ampersand replacement=\&]
 \mathrm{Diff}^{i}(P,F)\arrow[r,"{u\mapsto D\circ u}"]\&
 \mathrm{Diff}^{i+1}(P,G)\arrow[r,"{v\mapsto D'\circ v}"]\&
 \mathrm{Diff}^{i+2}(P,H)
\end{tikzcd}}
est exacte, et la suite
\diagram{\begin{tikzcd}[ampersand replacement=\&]
 \mathrm{Diff}^{i}(H,P)\arrow[r,"{u\mapsto u\circ D'}"]\&
 \mathrm{Diff}^{i+1}(G,P)\arrow[r,"{v\mapsto v\circ D}"]\&
 \mathrm{Diff}^{i+2}(F,P)
\end{tikzcd}}
est exacte.

% 2), Poincaré hypothesis
2) Supposons que $E$ sur $S$ soit loc libre \uline{stratifié}, et que le lemme de Poincaré formel soit vrai pour le complexe
\[
 E\otimes\Omega^*_{S/T}=\Omega^*_{S/T}(E)
\]
[p. ex. $S$ lisse sur $T$, et \unsure{caractéristique|caractéristiques} de $T$ \unsure{nulle|nulles}].

% leftward sequence, auxiliary Hom and its annotation
Alors \inserted{$\forall\,P$ qu. coh. et $i\geq 0$} la suite
\diagram{\scriptsize\begin{tikzcd}[ampersand replacement=\&,column sep=small,row sep=small]
 0\&\mathrm{Hom}(E,P)\arrow[l]\&\&
 \mathrm{Diff}^i(E,P)\arrow[ll,"{\deleted{\unsure{u\mapsto u(1)}}}"]
  \arrow[dl,"{\sim}"]\&
 \mathrm{Diff}^{i-1}(E\otimes\Omega^1_{S/T},P)\arrow[l,"{u\mapsto u\circ d_E^0}"]\&
 \cdots\arrow[l,"{u\mapsto u\circ d_E^1}"]\\
 \&\&\bigl[\mathrm{Hom}(P^i(E),P)\bigr]\arrow[ul,"{u\mapsto u\circ\unsure{\delta_i^{\mathrm{str}}|\delta_i^E|\delta_i}}"']\&\&\&
\end{tikzcd}}
\inserted{où $\delta_i^{\mathrm{str}}:E\longrightarrow P^i(E)$ défini par la stratification}
est exacte.

% NB at foot, continues
\uline{NB} $\operatorname{Ker}\bigl(\mathrm{Diff}^i(E,P)\longrightarrow\mathrm{Hom}(E,P)\bigr)$ est désigné par $\mathrm{Diff}^i_+(E,P)$, op. diff. ``à coeff. cts'' rel. à la stratification
% archive sheet number: 230

\scanpage{232}
% corollary and question
Donc on trouve

\uline{Corollaire}
\[
 \mathrm{Diff}^{(i)}_+(E,P)\simeq
 \mathrm{Diff}^{i-1}(E\otimes\Omega^1_{S/T},P)/
 \operatorname{Im}\mathrm{Diff}^{i-2}(E\otimes\Omega^{\unsure{2}}_{S/T},P)
\]
% right edge of the preceding denominator meets the scan boundary
Par suite, on a un hom naturel injectif \textcircled{?}
\[
 \mathrm{Diff}^{(i)}_+(E,P)\hookrightarrow
 \mathrm{Hom}_{\mathrm{add}}(\operatorname{Ker}d_E^1,P)
\]
\inserted{l'image est appelée : faisceau des op. diff. de $\operatorname{Ker}d_E^1$ dans $P$}

% universal target
Pour $i$ fixé, le $P$ ``universel'' pour les op. diff à coeff constants (d'ordre $\leq i$) sur $E$ est $\unsure{p|P}^{(i)}(E)/\underline{\mathcal O}_S\cdot\delta_i^{\deleted{\unsure{E}}}(E)$, où
\[
 \delta_i:E\longrightarrow p^i(E)
\]
est défini par la stratification.

% 3), submodule
3) Supposons maintenant donné un sous-Module $N$ de $E$, on s'intéresse aux hom ``différentiels'' $\operatorname{Ker}d_E^1\longrightarrow P$ qui s'annulent sur $d_E^0 N$, qui correspondent donc aux \deleted{\unsure{hom}} op. diff. $E\longrightarrow P$, à coeff. cts, et qui s'annulent sur $N$, donc aux op. diff. $M=E/N\longrightarrow P$ qui ``\deleted{\unsure{fon}} \deleted{sont nuls sur}'' deviennent à coeff. cts sur $E$. Ils correspondent \inserted{(pour l'ordre $\leq i$)} aux hom linéaires $p^i(M)\longrightarrow P$ qui s'annulent sur $\rho_i E$ ; où
\[
 \rho_i:E\longrightarrow p^i(M)
\]
% archive sheet number: 231

\scanpage{233}
% composite defining the map from the preceding leaf
est l'hom. composé [linéaire pour la structure \uline{gauche} sur $\unsure{P|p}^{i}(M)$]
\diagram{\begin{tikzcd}[ampersand replacement=\&]
 E\arrow[r,"{\delta_i}"]\&p^i(E)\arrow[r]\&p^i(M).
\end{tikzcd}}

% representing module and inverse system
Donc, pour $P$ variable, c'est représenté par
\[
 \Phi_i=p^i(M)/\underline{\mathcal O}_S\cdot\rho_i(E),
\]
les $\Phi_i$ \unsure{formant|forment} un syst. projectif strict, déterminé par la connaissance de $M$ et de
\[
 \rho_\infty:E\longrightarrow P^\infty(M)
\]
(linéaire pour la structure \uline{gauche} de $P^\infty(M)$), défini par les $\rho_i$. Attention, \unsure{ce|ici} $\rho_\infty$ satisfait à des \deleted{\unsure{deux}} conditions spéciales : son composé $E\longrightarrow M$ avec l'augmentation $P^\infty(M)\longrightarrow M$ est un épimorphisme, $+$ conditions d'intégrabilité---
% archive sheet number: 232

\scanpage{234}
% numbered heading and diagram of liftings
\mspage{\uline{1}}
\uline{Construction de la connexion canonique en cohomologie de De Rham}

\diagram{\begin{tikzcd}[ampersand replacement=\&,row sep=large,column sep=small]
 X_0\arrow[d,"{f_0}","{\substack{\text{schéma}\\\text{relatif}\\\text{formt. lisse}}}"']\&
 X'\arrow[d,"{f'}"]\&X''\arrow[dl,"{f''}"]\\
 T_0\arrow[r,phantom,"{\unsure{\longleftarrow|\subset|\hookrightarrow}}"]\&T\&
\end{tikzcd}}
% slanted; joined to the word lisse
\margin{au lieu de la lissité \inserted{de $X_0,X',X''$} il suffit que $X',X''$ soient loc. $T$-\unsure{isom|isomorphes|isom.} (sur $X_0$).}
immersion 1-nilpotente
\[
 T_0=V(\underline I)
\]
\[
 \begin{aligned}
 X'\times_T T_0&\xrightarrow{\sim}X_0\\
 X''\times_{\unsure{T|t}} T_0&\xrightarrow{\sim}X_0
 \end{aligned}
\]
$X',X''$ deux prolongements lisses de $X_0/T_0$ sur $T$.

% sheaves of germs; long repetition stroke in second line
\[
 \begin{aligned}
 \unsure{\mathscr P|\mathcal P}&\quad\text{faisceau sur esp }X_0\text{ des germes des }T\text{-isom }X'\longrightarrow X''\\
 \unsure{\mathscr G|\mathcal G}&\quad\rule{.42\linewidth}{.3pt}\qquad T\text{-autom }X''\longrightarrow X''
 \end{aligned}
\]
alors $\mathscr G\simeq\mathrm{Hom}(\Omega^1_{X_0/T_0},f_0^{\unsure{*|\alpha}}(I))=\unsure{\mathscr T|\mathcal T|\mathscr V}_{X_0/T_0}\otimes f_0^*(\underline I)$

$\mathscr P$ un \inserted{$\mathscr G$-torseur à droite} [faisceau principal homogène \inserted{à droite} sur $\mathscr G$. ]

% action on the de Rham complexes
\[
 \mathscr P\times\Omega^\bullet_{X'/T}\longrightarrow\Omega^\bullet_{X''/T}
 \qquad\text{i.e}
\]
\[
 \begin{aligned}
 \mathscr P&\xrightarrow{u}\underline{\text{Isom de complexes}}\,(\Omega^\bullet_{X'/T},\Omega^\bullet_{X''/T})\\
 &\subset \unsure{\mathscr Z|Z}^0\,\underline{\mathrm{Hom}}^\bullet(\Omega^\bullet_{X'/T},\Omega^\bullet_{X''/T})\\
 \mathscr G&\xrightarrow{v}\underline{\text{\deleted{\unsure{Isom}}Aut de complexes}}\,(\Omega^\bullet_{X''/T})
 \subset\mathscr Z^0\,\underline{\mathrm{Hom}}^\bullet(\Omega^\bullet_{X''/T},\Omega^\bullet_{X''/T}\unsure{)|.}\\
 \mathscr G&\xrightarrow{\unsure{k|h}}\underline{\mathrm{Hom}}^{(-1)}(\underline{\Omega}^\bullet_{X''/T},\Omega^\bullet_{X''/T})
 \quad\text{par \unsure{produit intérieur|produits intérieurs}}
 \end{aligned}
\]
On a
\[
 \left\{\begin{aligned}
 u(pg)&=u(p)v(g)&\text{si }p&\in\Gamma(U,\mathscr P)\\
 && g&\in\Gamma(U,\mathscr G)\\
 v(g)-\mathrm{id}&=dk(g)+k(g)d
 \end{aligned}\right.
\]
% archive sheet number: 233

\scanpage{235}
% page heading and abstract data
\mspage{\uline{2}.}
Plus généralement, sur un \unsure{espace|esp.}\ \deleted{top} \inserted{annelé} $X$, soient $K',K''$ deux complexes de Modules, et

$\unsure{\mathscr G|\mathcal G}$ un faisceau abélien

$\unsure{\mathscr P|\mathcal P}$ un $\mathscr G$-torseur à droite
\[
 \begin{aligned}
 \mathscr P&\xrightarrow{u}\deleted{\unsure{\mathbb R}}\unsure{\mathscr Z|Z}^0\underline{\mathrm{Hom}}^\bullet(K',K'')\\
 \mathscr G&\xrightarrow{v}\underline{\mathrm{Hom}}^\bullet(K'',K'')\\
 \mathscr G&\xrightarrow{\unsure{k|h}}\underline{\mathrm{Hom}}^{-1}(K^{\unsure{''|'}},K'')\\
 u(pg)&=u(p)v(g)\\
 v(g+g')&=v(g)+v(g')\\
 u(g)&=dk(g)+\unsure{h|k}(g)d
 \end{aligned}
\]

% intended map, in the derived category
On veut en conclure un hom dans la cat. dérivée $K'\longrightarrow K''$ d'où
\[
 \mathbb R\Gamma_X(K')\longrightarrow\mathbb R\Gamma_X(K'')
\]
(compatible avec la structure sur \deleted{\unsure{$\mathbb R\Gamma_X$}} $\Gamma_X(\underline{\mathcal O}_X)$....)

[et plus gén. en relativisant par rapport à un morphisme d'espaces annelés $X\unsure{\longrightarrow|-}T$ ...]

% the extension and its associated complex
On introduit l'extension $\overline{\mathscr G}$ de $\mathbb Z$ par $\mathscr G$
\diagram{\begin{tikzcd}[ampersand replacement=\&]
 0\arrow[r]\&\mathscr G\arrow[r,"i"]\&\overline{\mathscr G}\arrow[r,"j"]\&\mathbb Z\arrow[r]\&0
\end{tikzcd}}
\unsure{défini|définie} \unsure{au travers|en termes|au dessus} de $\mathscr P$ par la relation
\[
 j^{-1}(1)\simeq\mathscr P\qquad
 (\text{isom de }\deleted{\unsure{\text{groupes}|\mathscr G\text{-is}}}\uline{\mathscr G\text{-torseurs}}\ldots).
\]
Soit
\[
 \mathscr H=\underline{\mathrm{Hom}}^\bullet(K',K'')
\]
et soit $\mathcal L$ le complexe
\diagram{\begin{tikzcd}[ampersand replacement=\&,column sep=small]
 0\arrow[r]\&0\arrow[r]\&\mathscr G^{\unsure{=|''}\mathcal L^0}\arrow[r]\&
 \overline{\mathscr G}^{=\mathcal L^1}\arrow[r]\&0\arrow[r]\&0
\end{tikzcd}}
% archive sheet number: 234

\scanpage{236}
% morphism of complexes
On trouve alors un homomorphisme
\[
 \varphi:\mathcal L^\bullet\longrightarrow\mathscr H^\bullet
\]
par
\diagram{\begin{tikzcd}[ampersand replacement=\&,column sep=small,row sep=large]
 \cdots\arrow[r]\&0\arrow[r]\&\unsure{\mathscr G|\mathcal G}\arrow[r,"i"]\arrow[d,"{\deleted{\unsure{k}}\varphi^{-1}}"]\&
 \overline{\mathscr G}\arrow[r]\arrow[d,"{\varphi^0}"]\&0\arrow[r]\&\cdots\\
 \&\mathscr H^{-2}\arrow[r]\&\mathscr H^{-1}\arrow[r,"{\delta^{\unsure{-1}}}"]\&
 \mathscr H^0\arrow[r]\&\mathscr H^1\arrow[r]\&\cdots
\end{tikzcd}}
défini par les conditions
\[
 \varphi^{-1}=\unsure{k|h}\qquad \varphi^0|\unsure{\mathscr P|\mathcal P}=u
\]
et la commutativité du carré

% explanation of a homomorphism out of the extension
[\deleted{Du fait que} \inserted{Un} hom de $\overline{\mathscr G}$ dans un un faisceau $F$ $=$ \{ Hom $v$ de $\mathscr G^{\unsure{\cdot\cdot}}$ dans $F$, et hom $u$ de $\mathscr P$ dans $F$, avec condition $u(pg)=u(p)v(g)$\}

Donc prenant \deleted{$\delta\varphi^{-1}$} $F=\mathscr H^0$, et $u,v$ comme dessus, on trouve $\varphi^0$, avec $\varphi^0 i=u$. Mais $u=\delta\deleted{\unsure{\varphi^{-1}|\varphi}}k=\delta\varphi^{-1}$ ... ]

% tensoring, preliminary composite and boxed result
Tensorisant avec $\underline{\mathcal O}_X$, on trouve
\diagram{\scriptsize\begin{tikzcd}[ampersand replacement=\&,column sep=small]
 \underline{\mathcal O}_X\arrow[r,"{\text{dans cat.}}","{\text{dérivée}}"']\&
 \deleted{\unsure{\mathcal L^0|\mathscr L}}\inserted{\unsure{\mathcal L^\bullet|\mathscr L}}\arrow[r]\&\mathscr H\arrow[r]\&
 \mathbb R\underline{\mathrm{Hom}}(K',K'')
\end{tikzcd}}
\deleted{On trouve donc,} donc
\[
 \boxed{\underline{\mathcal O}_X\unsure{\xrightarrow{\sim}|\longrightarrow}
 \mathbb R\underline{\mathrm{Hom}}(K',K'')\quad\text{ds cat. dérivée}}
\]
Donc
\[
 \deleted{\mathbb R\Gamma_X(\unsure{\mathcal L})}\simeq
 \mathbb R\Gamma_X(\deleted{\unsure{\mathbb Z}}\underline{\mathcal O}_X)\longrightarrow
 \deleted{\mathbb R\Gamma_X\unsure{\mathscr H}}\longrightarrow\mathbb R\mathrm{Hom}(K',K'')
\]
d'où
\[
 \boxed{\mathbb R\Gamma_X(\deleted{\unsure{\mathbb Z}}\underline{\mathcal O}_X)\longrightarrow\mathbb R\mathrm{Hom}(K',K'')}
\]
Par suite, \deleted{faisant} passant au $\mathbb R\Gamma^0_X$, on trouve un \unsure{él\textsuperscript{mt}|élémt|élément} de $\mathbb R^0\mathrm{Hom}(K',K'')$, i.e. un hom de la cat\unsure{e|\textsuperscript{e}|é|.} dérivée
% archive sheet number: 235
